\documentclass[10pt,a4paper]{amsart}
\usepackage[margin=19mm]{geometry}
\usepackage{amsmath,amssymb,mathtools}
\usepackage[scr=rsfs,cal=euler]{mathalfa}
\usepackage{xfrac}
\usepackage[unicode,hidelinks,bookmarks=false]{hyperref}
\newcommand{\ie}{i.e.\ }
\newcommand{\cf}{cf.\ }
\newcommand{\resp}{respectively\ }
\newcommand{\Subsection}{\subsection}
\providecommand{\nobreakdash}{\nobreak\hbox{-}}
\setlength{\emergencystretch}{2em}
\multlinegap0pt
\usepackage{tikz}
\usepackage{xcolor}
\usepackage{booktabs}
\usepackage{subfig}
\usepackage{amssymb}
\usepackage[shortlabels]{enumitem}
\newtheorem{theorem}{Theorem}
\title{On the Structure of the Optimal Detector for Sub-THz Multi-Hop Relays with Unknown Prior: Over-the-Air Diffusion}
\author{Ozgur Ercetin, Mohaned Chraiti}
\date{Source: arXiv:2601.01194v1 (CC BY 4.0)}
\begin{document}
\maketitle

\noindent\emph{Source and licence.} On the Structure of the Optimal Detector for Sub-THz Multi-Hop Relays with Unknown Prior: Over-the-Air Diffusion. Version arXiv:2601.01194v1 (CC BY 4.0), distributed by arXiv under the Creative Commons Attribution 4.0 International licence, http://creativecommons.org/licenses/by/4.0/, as recorded in the arXiv licence metadata for this exact version and shown on the abstract page https://arxiv.org/abs/2601.01194v1. Full licence notice: \texttt{licenses/arxiv-2601.01194-CC-BY-4.0.txt}.\par\smallskip

\noindent\textit{Editorial scope.} Counted unit: the article's theorem thm:equalization (source0.tex lines 544--573). The article's complete original proof of it is delivered with it (source0.tex lines 575--628, one \texttt{proof} environment of 54 lines). No mathematics is written, repaired or re-derived: the statement and the proof are the article's own lines, in the article's own notation, and no retained line is substituted or re-flowed.  No further source-local context is needed: every cross-reference of every retained line resolves inside the two delivered spans.  Nothing was omitted from any retained span, so the package carries no omission record.  No retained line contains a citation, so the wrapper carries no bibliography and no bibliography entry.  No retained line contains a figure environment, an image-inclusion command or drawing code, so no image is delivered and the package declares no asset dependency.  Editorial formatting only, in addition: an AMS \texttt{amsart} preamble that restates the article's own declarations and macros used by the retained lines, namely the environments \texttt{theorem} on separate counters, together with the standard packages those lines need.  The retained environments print in this document as Theorem 1 in order of appearance, whereas the article prints the same items with the numbering of its own full document.  The complete native source of the article as distributed in the arXiv e-print for this exact version is bundled unchanged as \texttt{sources/192060/source0.tex}.
\begin{theorem}[Optimal Power Allocation Across Relays]\label{thm:equalization}
Consider the recursion with $\alpha_t^2 = 1-\beta_t$ and
\[
\beta_t = \frac{1}{1+\mathrm{SNR}_{\mathrm{in},t}}
= \frac{1}{1+c_t P_t}, 
\]
where $P_t\in[0,P_{t,\max}]$ is the transmit power of relay $t$, 
and $c_t=\frac{|h_t|^2}{\sigma_t^2}>0$ is the channel-to-noise ratio parameter at hop $t$. 
Then
\[
\bar\alpha_H^2=\prod_{t=0}^{T-1}(1-\beta_t)
=\prod_{t=0}^{T-1}\frac{c_tP_t}{1+c_tP_t}.
\]

Given a total power budget $P_{\mathrm{tot}}>0$, the optimal allocation is obtained by solving
\begin{align}\label{eq:opt}
&\max_{\{P_t\}}\ \sum_{t=0}^{T-1}\log\frac{c_tP_t}{1+c_tP_t}\\
&\quad \text{s.t.}\quad
\sum_{t=0}^{T-1} P_t \le P_{\mathrm{tot}},\qquad 
0 \le P_t \le P_{t,\max}.
\nonumber
\end{align}
This optimization is strictly concave and admits a unique solution $\{P_t^\star\}$. 
For each interior variable $0<P_t^\star<P_{t,\max}$, the Karush--Kuhn--Tucker (KKT) 
conditions reduce to the equal-marginal rule
\[
\frac{1}{P_t^\star(1+c_tP_t^\star)}=\mu,
\]
for some multiplier $\mu>0$ chosen so that the total budget is met. 
\end{theorem}

\begin{proof}
Define
\[
g_t(P) = \log\frac{c_tP}{1+c_tP}, \qquad P>0,\ c_t>0.
\]
We first verify concavity. The first and second derivatives are
\[
g_t'(P) = \frac{1}{P(1+c_tP)} > 0, \qquad
g_t''(P) = -\frac{1+2c_tP}{P^2(1+c_tP)^2} < 0.
\]
Hence $g_t(P)$ is strictly increasing and strictly concave on $(0,\infty)$. 
It follows that the sum $\sum_t g_t(P_t)$ is strictly concave and that 
the feasible set in \eqref{eq:opt} is convex. Therefore, the problem admits 
a unique optimizer.

Next, consider the Lagrangian
\[
\begin{aligned}
\mathcal{L} &= \sum_{t=0}^{T-1} g_t(P_t) 
- \mu\Bigl(\sum_{t=0}^{T-1} P_t - P_{\mathrm{tot}}\Bigr)\\
&\qquad+ \sum_{t=0}^{T-1} \nu_t P_t 
+ \sum_{t=0}^{T-1} \omega_t (P_{t,\max}-P_t),
\end{aligned}
\]
where $\mu\ge0$ is the multiplier for the total power constraint, 
and $\nu_t,\omega_t\ge0$ are multipliers for the box constraints.

The KKT stationarity condition yields
\[
g_t'(P_t^\star) = \mu - \nu_t + \omega_t.
\]
For any interior solution $0<P_t^\star<P_{t,\max}$, the complementary slackness 
conditions imply $\nu_t=\omega_t=0$, and therefore
\[
g_t'(P_t^\star) = \mu.
\]
Substituting the explicit form of $g_t'(P)$ gives
\[
\frac{1}{P_t^\star(1+c_tP_t^\star)} = \mu,
\]
which can be rearranged as
\[
P_t^\star(1+c_tP_t^\star) = \frac{1}{\mu}.
\]

This quadratic equation in $P_t^\star$ has the positive solution
\[
P_t^\star = \frac{\sqrt{1+\tfrac{4c_t}{\mu}}-1}{2c_t}.
\]
Boundary cases $P_t^\star=0$ or $P_t^\star=P_{t,\max}$ follow from the 
corresponding slackness conditions. Finally, since $g_t'(P)>0$, the 
budget constraint is tight unless prevented by the upper bounds 
$P_{t,\max}$. 
\end{proof}

\end{document}
